\documentclass[../../main.tex]{subfiles}

\begin{document}

\section{Lebesgue 可测集的 $\sigma$ 代数}\label{section:ch02-可测集}

外测度具有四个优点: (i) 对实数的所有集合定义; (ii) 区间的外测度是它的长度; (iii) 外测度是可数次可加的; (iv) 外测度是平移不变的. 但外测度不是可数可加的. 事实上, 它甚至不是有限可加的（见\cref{theorem:ch02-不可加}）: 存在不相交的集合 $A$ 和 $B$ 使得
\begin{align}
m^{*}(A\cup B)<m^{*}(A)+m^{*}(B).\label{align:ch02-严格不等}
\end{align}

为改善这个基本的缺陷, 我们支持称为 Lebesgue 可测集的集合的 $\sigma$ 代数, 它包含所有区间与所有开集且具有性质: 外测度这一集函数限制在可测集族上是可数可加的. 有若干个方式定义一个集合是可测的意味着什么. 我们遵循 Constantin Carath\'eodory 的方式.

\begin{definition}[Carath\'eodory 可测]\label{definition:ch02-Caratheodory}
集合 $E$ 称为是可测的, 若对任意集合 $A$,
\begin{align*}
m^{*}(A)=m^{*}(A\cap E)+m^{*}(A\cap E^{c}).
\end{align*}
\end{definition}

我们立即看到可测集具有的一个优越性, 即若集合 $A$ 与 $B$ 中的一个是可测的, 严格的不等式\eqref{align:ch02-严格不等}不可能发生. 事实上, 若 $A$ 是可测的而 $B$ 是任何与 $A$ 不交的集合, 则
\begin{align*}
m^{*}(A\cup B)=m^{*}([A\cup B]\cap A)+m^{*}([A\cup B]\cap A^{c})=m^{*}(A)+m^{*}(B).
\end{align*}

根据\cref{proposition:ch02-外测度次可加}, 由于外测度是有限次可加的而 $A=[A\cap E]\cup[A\cap E^{c}]$, 我们总是有
\begin{align*}
m^{*}(A)\leqslant m^{*}(A\cap E)+m^{*}(A\cap E^{c}).
\end{align*}
因此 $E$ 是可测的当且仅当对每个集合 $A$ 我们有
\begin{align}
m^{*}(A)\geqslant m^{*}(A\cap E)+m^{*}(A\cap E^{c}).\label{align:ch02-可测不等式}
\end{align}
若 $m^{*}(A)=\infty$, 该不等式平凡成立. 因此, 仅须对具有有限外测度的集合 $A$ 证明\eqref{align:ch02-可测不等式}.

观察到可测性的定义关于 $E$ 和 $E^{c}$ 是对称的, 因此一个集合是可测的当且仅当它的补是可测的. 显然空集 $\varnothing$ 和所有实数构成的集合 $\mathbf{R}$ 是可测的.

\begin{proposition}\label{proposition:ch02-零外测度可测}
任何外测度为零的集合是可测的. 特别地, 任何可数集是可测的.
\end{proposition}

\begin{proof}
令集合 $E$ 的外测度为零. 令 $A$ 为任意集合. 由于
\begin{align*}
A\cap E\subseteq E\quad\text{且}\quad A\cap E^{c}\subseteq A,
\end{align*}
根据外测度的单调性,
\begin{align*}
m^{*}(A\cap E)\leqslant m^{*}(E)=0\quad\text{且}\quad m^{*}(A\cap E^{c})\leqslant m^{*}(A).
\end{align*}
因此
\begin{align*}
m^{*}(A)\geqslant m^{*}(A\cap E^{c})=0+m^{*}(A\cap E^{c})=m^{*}(A\cap E)+m^{*}(A\cap E^{c}).
\end{align*}
从而 $E$ 是可测的.
\end{proof}

\begin{proposition}\label{proposition:ch02-有限并可测}
可测集的有限族的并是可测的.
\end{proposition}

\begin{proof}
第一步, 我们证明两个可测集 $E_{1}$ 和 $E_{2}$ 的并是可测的. 令 $A$ 为任意集合. 首先用 $E_{1}$ 的可测性, 接着用 $E_{2}$ 的可测性, 我们有
\begin{align*}
m^{*}(A)&=m^{*}(A\cap E_{1})+m^{*}(A\cap E_{1}^{c})\\
&=m^{*}(A\cap E_{1})+m^{*}([A\cap E_{1}^{c}]\cap E_{2})+m^{*}([A\cap E_{1}^{c}]\cap E_{2}^{c}).
\end{align*}
以下集合等式成立:
\begin{align*}
[A\cap E_{1}^{c}]\cap E_{2}^{c}=A\cap[E_{1}\cup E_{2}]^{c}
\end{align*}
与
\begin{align*}
[A\cap E_{1}]\cup[A\cap E_{1}^{c}\cap E_{2}]=A\cap[E_{1}\cap E_{2}].
\end{align*}
我们从这些等式和外测度的有限次可加性推出
\begin{align*}
m^{*}(A)&=m^{*}(A\cap E_{1})+m^{*}([A\cap E_{1}^{c}]\cap E_{2})+m^{*}([A\cap E_{1}^{c}]\cap E_{2}^{c})\\
&=m^{*}(A\cap E_{1})+m^{*}([A\cap E_{1}^{c}]\cap E_{2})+m^{*}(A\cap[E_{1}\cup E_{2}]^{c})\\
&\geqslant m^{*}(A\cap[E_{1}\cup E_{2}])+m^{*}(A\cap[E_{1}\cup E_{2}]^{c}).
\end{align*}
因此 $E_{1}\cup E_{2}$ 是可测的.

现在令 $\{E_{k}\}_{k=1}^{n}$ 为任意可测集的有限族. 对于一般的 $n$, 我们用归纳法证明并集 $\bigcup_{k=1}^{n}E_{k}$ 的可测性. 对 $n=1$ 这是平凡的. 假定它对 $n-1$ 成立. 因此, 由于
\begin{align*}
\bigcup_{k=1}^{n}E_{k}=\left[\bigcup_{k=1}^{n-1}E_{k}\right]\cup E_{n},
\end{align*}
且我们已证明了两个可测集的并的可测性, 集合 $\bigcup_{k=1}^{n}E_{k}$ 是可测的.
\end{proof}

\begin{proposition}\label{proposition:ch02-有限不交可加}
令 $A$ 为任意集合而 $\{E_{k}\}_{k=1}^{n}$ 是可测集的有限不交族. 则
\begin{align*}
m^{*}\left(A\cap\left[\bigcup_{k=1}^{n}E_{k}\right]\right)=\sum_{k=1}^{n}m^{*}(A\cap E_{k}).
\end{align*}
特别地,
\begin{align*}
m^{*}\left(\bigcup_{k=1}^{n}E_{k}\right)=\sum_{k=1}^{n}m^{*}(E_{k}).
\end{align*}
\end{proposition}

\begin{proof}
证明通过对 $n$ 归纳进行. 对 $n=1$ 显然结论成立. 假设命题对 $n-1$ 成立. 由于族 $\{E_{k}\}_{k=1}^{n}$ 是互不相交的,
\begin{align*}
A\cap\left[\bigcup_{k=1}^{n}E_{k}\right]\cap E_{n}=A\cap E_{n}
\end{align*}
且
\begin{align*}
A\cap\left[\bigcup_{k=1}^{n}E_{k}\right]\cap E_{n}^{c}=A\cap\left[\bigcup_{k=1}^{n-1}E_{k}\right].
\end{align*}
因此, 根据 $E_{n}$ 的可测性以及归纳假设,
\begin{align*}
m^{*}\left(A\cap\left[\bigcup_{k=1}^{n}E_{k}\right]\right)&=m^{*}(A\cap E_{n})+m^{*}\left(A\cap\left[\bigcup_{k=1}^{n-1}E_{k}\right]\right)\\
&=m^{*}(A\cap E_{n})+\sum_{k=1}^{n-1}m^{*}(A\cap E_{k})=\sum_{k=1}^{n}m^{*}(A\cap E_{k}).
\end{align*}
\end{proof}

$\mathbf{R}$ 的子集族称为是一个代数, 若它包含 $\mathbf{R}$ 且它关于补与有限并封闭; 根据 De Morgan 等式, 这样的族也关于有限交封闭. 我们从\cref{proposition:ch02-有限并可测}以及可测集的补集的可测性推出可测集族是一个代数. 观察到可测集的可数族的并也是可测集的可数不交族的并是有用的. 事实上, 令 $\{A_{k}\}_{k=1}^{\infty}$ 为可测集的可数族. 定义 $A_{1}'=A_{1}$ 且对每个 $k\geqslant 2$, 定义
\begin{align*}
A_{k}'=A_{k}\sim\bigcup_{i=1}^{k-1}A_{i}.
\end{align*}
由于可测集族是一个代数, $\{A_{k}'\}_{k=1}^{\infty}$ 是可测集的互不相交族, 其并与 $\{A_{k}\}_{k=1}^{\infty}$ 的并相同.

\begin{proposition}\label{proposition:ch02-可数并可测}
可测集的可数族的并是可测的.
\end{proposition}

\begin{proof}
令 $E$ 为可测集的可数族的并. 如同我们上面观察到的, 存在可测集的可数不交族 $\{E_{k}\}_{k=1}^{\infty}$ 使得 $E=\bigcup_{k=1}^{\infty}E_{k}$. 令 $A$ 为任意集合. 令 $n$ 为自然数. 定义 $F_{n}=\bigcup_{k=1}^{n}E_{k}$. 由于 $F_{n}$ 是可测的且 $F_{n}^{c}\supseteq E^{c}$,
\begin{align*}
m^{*}(A)=m^{*}(A\cap F_{n})+m^{*}(A\cap F_{n}^{c})\geqslant m^{*}(A\cap F_{n})+m^{*}(A\cap E^{c}).
\end{align*}
根据\cref{proposition:ch02-有限不交可加},
\begin{align*}
m^{*}(A\cap F_{n})=\sum_{k=1}^{n}m^{*}(A\cap E_{k}).
\end{align*}
因此
\begin{align*}
m^{*}(A)\geqslant\sum_{k=1}^{n}m^{*}(A\cap E_{k})+m^{*}(A\cap E^{c}).
\end{align*}
该不等式左边与 $n$ 无关. 因此
\begin{align*}
m^{*}(A)\geqslant\sum_{k=1}^{\infty}m^{*}(A\cap E_{k})+m^{*}(A\cap E^{c}).
\end{align*}
因此, 根据外测度的可数次可加性,
\begin{align*}
m^{*}(A)\geqslant m^{*}(A\cap E)+m^{*}(A\cap E^{c}).
\end{align*}
因此 $E$ 是可测的.
\end{proof}

$\mathbf{R}$ 的子集族称为 $\sigma$ 代数, 若它包含 $\mathbf{R}$ 且关于补与可数并封闭; 根据 De Morgan 等式, 这样的族关于可数交也封闭. 前一个命题告诉我们可测集族是一个 $\sigma$ 代数.

\begin{proposition}\label{proposition:ch02-区间可测}
每个区间是可测的.
\end{proposition}

\begin{proof}
如同我们上面观察到的, 可测集族是一个 $\sigma$ 代数. 因此为证明每个区间是可测的, 仅须证明每个形如 $(a,\infty)$ 的区间是可测的（见习题 11）. 考虑这样的一个区间. 令 $A$ 为任意集合. 我们假定 $a$ 不属于 $A$. 否则, 用 $A\sim\{a\}$ 代替 $A$, 保持外测度不变. 我们必须证明
\begin{align}
m^{*}(A_{1})+m^{*}(A_{2})\leqslant m^{*}(A),\label{align:ch02-区间可测证}
\end{align}
其中
\begin{align*}
A_{1}=A\cap(-\infty,a)\quad\text{而}\quad A_{2}=A\cap(a,\infty).
\end{align*}

根据 $m^{*}(A)$ 作为下确界的定义, 为证明\eqref{align:ch02-区间可测证}, 必要且充分的是证明对任何覆盖 $A$ 的开有界区间的可数族 $\{I_{k}\}_{k=1}^{\infty}$,
\begin{align}
m^{*}(A_{1})+m^{*}(A_{2})\leqslant\sum_{k=1}^{\infty}\ell(I_{k}).\label{align:ch02-覆盖拆分}
\end{align}
事实上, 对这样的一个覆盖, 对每个指标 $k$, 定义
\begin{align*}
I_{k}'=I_{k}\cap(-\infty,a)\quad\text{与}\quad I_{k}''=I_{k}\cap(a,\infty).
\end{align*}
则 $I_{k}'$ 和 $I_{k}''$ 是区间且
\begin{align*}
\ell(I_{k})=\ell(I_{k}')+\ell(I_{k}'').
\end{align*}
由于 $\{I_{k}'\}_{k=1}^{\infty}$ 和 $\{I_{k}''\}_{k=1}^{\infty}$ 分别是覆盖 $A_{1}$ 和 $A_{2}$ 的开有界区间的可数族, 根据外测度的定义,
\begin{align*}
m^{*}(A_{1})\leqslant\sum_{k=1}^{\infty}\ell(I_{k}'),\qquad m^{*}(A_{2})\leqslant\sum_{k=1}^{\infty}\ell(I_{k}'').
\end{align*}
因此
\begin{align*}
m^{*}(A_{1})+m^{*}(A_{2})\leqslant\sum_{k=1}^{\infty}\ell(I_{k}')+\sum_{k=1}^{\infty}\ell(I_{k}'')=\sum_{k=1}^{\infty}\left[\ell(I_{k}')+\ell(I_{k}'')\right]=\sum_{k=1}^{\infty}\ell(I_{k}).
\end{align*}
因此\eqref{align:ch02-覆盖拆分}成立且证明完成.
\end{proof}

每个开集是可数个不交开区间的并. 我们因此从前两个命题推出每个开集是可测的. 每个闭集是开集的补, 且因此每个闭集是可测的. 回忆一个实数的集合称为是一个 $G_{\delta}$ 集, 若它是可数个开集族的交; 而称为是一个 $F_{\sigma}$ 集, 若它是可数个闭集族的并. 我们从\cref{proposition:ch02-可数并可测}推断出每个 $G_{\delta}$ 集与每个 $F_{\sigma}$ 集是可测的.

所有包含开集的 $\mathbf{R}$ 的子集的 $\sigma$ 代数的交是一个 $\sigma$ 代数, 它称为 Borel $\sigma$ 代数, 该族的成员称为 Borel 集. Borel $\sigma$ 代数包含在每个包含所有开集的 $\sigma$ 代数中. 因此, 由于可测集是包含所有开集的 $\sigma$ 代数, 每个 Borel 集是可测的. 我们已建立了以下定理.

\begin{theorem}\label{theorem:ch02-可测Borel}
可测集族 $\mathcal{M}$ 是一个包含 Borel 集的 $\sigma$ 代数 $\mathcal{B}$ 的 $\sigma$ 代数. 每个区间、每个开集、每个闭集、每个 $G_{\delta}$ 集和每个 $F_{\sigma}$ 集是可测的.
\end{theorem}

\begin{proposition}\label{proposition:ch02-可测平移}
可测集的平移是可测的.
\end{proposition}

\begin{proof}
令 $E$ 为可测集. 令 $A$ 为任意集合, 而 $y$ 是实数. 根据 $E$ 的可测性和外测度的平移不变性,
\begin{align*}
m^{*}(A)&=m^{*}(A-y)=m^{*}([A-y]\cap E)+m^{*}([A-y]\cap E^{c})\\
&=m^{*}(A\cap[E+y])+m^{*}(A\cap[E+y]^{c}).
\end{align*}
因此 $E+y$ 是可测的.
\end{proof}

\begin{exercise}\label{exercise:ch02-11}
证明: 若 $\mathbf{R}$ 的子集的 $\sigma$ 代数包含形如 $(a,\infty)$ 的区间, 则它包含所有区间.
\end{exercise}

\begin{exercise}\label{exercise:ch02-12}
证明每个区间是一个 Borel 集.
\end{exercise}

\begin{exercise}\label{exercise:ch02-13}
证明: (i) $F_{\sigma}$ 集的平移也是 $F_{\sigma}$; (ii) $G_{\delta}$ 集的平移也是 $G_{\delta}$; (iii) 测度为零的集合的平移测度也为零.
\end{exercise}

\begin{exercise}\label{exercise:ch02-14}
证明: 若一个集合 $E$ 具有正的外测度, 则存在 $E$ 的有界子集也有正的外测度.
\end{exercise}

\begin{exercise}\label{exercise:ch02-15}
证明: 若 $E$ 有有限测度且 $\varepsilon>0$, 则 $E$ 是有限个可测集的不交并, 这些可测集的每个集合的测度至多为 $\varepsilon$.
\end{exercise}

\end{document}